% TOP_PROBLEM: {"id": 3238, "title": "Graphs with a forbidden induced tree are chi-bounded", "category_id": 3, "category": {"id": 3, "name": "graph_theory", "display_name": "Graph Theory", "description": "Problems involving graphs, networks, and their properties.", "slug": "graph-theory", "order_index": 3, "created_at": "2026-07-31T15:26:25.671Z"}, "status": "open", "problem_number": "OPG-36936", "set_id": 12}
% FIRST_POSTED: 2026-10-01
% ENTRY_KIND: statement_only
% UnsolvedMath ID 3238; source catalogue code OPG-36936
% !TeX program = lualatex
% Definition and sources only; this file is not an LLM solution attempt.
\documentclass[11pt,a4paper]{article}
\usepackage[margin=25mm]{geometry}
\usepackage{fontspec}
\setmainfont{lmroman10-regular.otf}[BoldFont=lmroman10-bold.otf,
 ItalicFont=lmroman10-italic.otf,BoldItalicFont=lmroman10-bolditalic.otf]
\usepackage{amsmath,amssymb,mathtools}
\usepackage{unicode-math}
\setmathfont{latinmodern-math.otf}
\AtBeginDocument{\let\setminus\smallsetminus}
\usepackage{xurl,hyperref}
\hypersetup{colorlinks=true,urlcolor=blue}
\setcounter{secnumdepth}{0}
\setlength{\parindent}{0pt}
\setlength{\parskip}{5pt}
\setlength{\emergencystretch}{3em}
\begin{document}
\section*{Graphs with a forbidden induced tree are chi-bounded}\label{3238}
{\small UnsolvedMath ID 3238 · OPG-36936 · Graph Theory · OpenGarden}
\subsection{Definitions and mathematical statement}
All graphs here are finite, simple, and undirected. For $U\subseteq V(G)$,
the induced subgraph $G[U]$ contains exactly the edges of $G$ with both ends
in $U$. A tree is a connected graph with no cycle. Fix a finite tree $T$;
a graph is $T$-free in the present sense when none of its induced subgraphs
is isomorphic to $T$.

Write $\chi(G)$ for the minimum number of colours in a proper vertex
colouring and $\omega(G)$ for the largest size of a clique, that is, a set
of pairwise adjacent vertices. Set both parameters to zero for the empty
graph. A graph class $\mathcal C$ is $\chi$-bounded if some function
$f:\mathbb N_0\to\mathbb N_0$ satisfies
$\chi(G)\le f(\omega(G))$ for every $G\in\mathcal C$.

The conjecture asks whether, for every fixed tree $T$, there exists a
function $f_T$ such that
\[
  G\text{ has no induced copy of }T
       \quad\Longrightarrow\quad\chi(G)\le f_T(\omega(G))
\]
for every finite graph $G$. The function may depend on $T$ but not on
$|V(G)|$. Equivalently, for each tree $T$ and each integer $w\ge1$ there
should be a finite bound $b(T,w)$ on the chromatic numbers of all
$T$-free graphs with clique number at most $w$. No polynomial or effective
bound is required. Replacing an induced copy by an arbitrary subgraph
would change the problem: extra edges among the selected vertices are
forbidden in the stated target.
\subsection{Short English statement}
Forbidding a fixed tree as an induced subgraph should bound the chromatic number whenever the clique number is bounded.
\subsection{Sources}
\begin{itemize}
\item[S1] ULAM.AI and the UnsolvedMath Contributors, supplied problems.json, record 3238 (OPG-36936), OpenGarden collection. Catalogue identity and source transcription; CC BY 4.0.
\url{https://huggingface.co/datasets/ulamai/UnsolvedMath}.
\item[S2] Open Problem Garden, Graphs with a forbidden induced tree are chi-bounded. Original problem and discussion; the mathematical formulation below is expanded from the supplied source record.
\url{http://www.openproblemgarden.org/op/graphs_with_a_forbidden_induced_tree_are_chi_bounded}.
\end{itemize}
\end{document}
